\begin{thebibliography}{}
\label{sec:Literatur}
\bibitem[1]{Karm2011} Karmasin, Matthias/Ribing, Rainer: Die Gestaltung wissenschaftlicher Arbeiten/Ein Leitfaden f�r Seminararbeiten, Bachelor-, Master- und Magisterarbeiten sowie Dissertationen, Wien, 2011, 6. Auflage
\bibitem[2]{Balz2008} Balzert, Helmut/Sch�fer, Christian/Schr�der, Marion: Wissenschaftliches Arbeiten/Wissenschaft, Quellen, Artefakte, Organisation, Pr�sentation, Herdecke/Witten, 2008, 2. Auflage
\bibitem[3]{Kipp2012} HSHL, Skript Wissenschaftliches Arbeiten, Referent: M.A. MLIS Guido Kippelt, Stand 10.05.2012
\bibitem[4]{Hering2009} Hering, Lutz/Hering, Heike: Technische Berichte, Verst�ndlich gliedern, gut gestalten,
�berzeugend vortragen, Vieweg+Teubner | GWV Fachverlage GmbH, Wiesbaden 2009, 6., aktualisierte und erweiterte Auflage
\bibitem[5]{Mess2012} HSHL, Formelsammlung Mess- und Regelungstechnik, WS 2011/12, Prof. Dr.-Ing. Ulrich Schneider
\end{thebibliography}

% Fu�noten bei Zitaten verzichten!
% Quellen in der Reihenfolge, wie sie im Text erscheinen!
% Im Text in normaler Gr��e Nummer in eckiger Klammer einf�gen [1]
% Im Literaturverzeichnis steht [1] und die Erkl�rung
% Literaturverzeichnis MUSS im Inhaltsverzeichnis erscheinen!